Conventions as in the statement: the store sign of Fact \claimref{synthesis.imported-gtg-dictionary}, $\hbar = c = G = 1$, $t = 2\pi T$. The check computes every curvature below from the gauge field itself, as exact Taylor jets of $\hb^{-1}$ about a rational point, the torsion-free $\Omega$ solved from them and the field strength read off.

\emph{The Clausius defect in this sign.} By Theorem \claimref{synthesis.einstein-equations-from-c-loc-and-clausius}(2), on a thin pencil about the generator through $p$ the defect $\delta Q - T_{\chi}\,\eta\,\delta A$ is $\tfrac{\kappa}{2}\bigl(T(k,k) - \tfrac{\eta}{2\pi}\Ric^{\circ}(k,k)\bigr)\lambda_{0}^{2}$ at leading order, with $\Ric^{\circ} = -\Ric$, the store's Ricci map being $\mathcal{R}(b)$. Hence the defect is $\tfrac{\kappa}{4\pi}\bigl(t(k,k) + \eta\,k\cdot\mathcal{R}(k)\bigr)\lambda_{0}^{2}$, and for a density constant on a connected field (H1) holds on every patch exactly when $t(k,k) + \eta\,k\cdot\mathcal{R}(k) = 0$ for every null $k$, that is (polarising over null directions) when the field equations hold with $G_{e} = 1/(4\eta)$ and some $\Lambda$, as in that theorem.

\emph{(i).} Let $U$ be an admissible field and $p_{n} \in U$ with local data converging to calibrated data $d_{*}$. Since $\eta = f(\text{local data})$ is $C^{1}$ on $U$ and (H1) holds on every patch, Theorem \claimref{synthesis.c-loc-entropy-clause-curvature-independent-density}(iv) gives $k\cdot D\eta = 0$ for every null $k$, hence $\eta$ is constant on the connected $U$. Continuity of $f$ gives $\eta = \lim f(\text{data}(p_{n})) = f(d_{*})$. For de Sitter data $f(d_{*}) = \tfrac14$: the de Sitter field's local data are the same at every point of a cut of its cosmological horizon, so the integrated density over a cut of area $4\pi N^{2}$ is $f(d_{*})\cdot 4\pi N^{2}$, and the stored de Sitter entropy $\pi N^{2}$ under the identification of Theorem \claimref{synthesis.c-loc-entropy-clause-curvature-independent-density}(ii) gives $f(d_{*}) = \tfrac14$. For flat data $f(0) = \lim_{\Lambda\to 0} f(\tfrac{\Lambda}{3}B, 0) = \tfrac14$ by continuity. So $\eta = \tfrac14$ on $U$, and by the defect identity the field equations hold with $G_{e} = 1/(4\eta) = 1$.

\emph{(ii).} For $\hb^{-1}(e_{i}) = a(t)\,e_{i}$ the coframe depends on $t$ alone and the curvature at a point depends on $(a, \dot a, \ddot a)$ there; with $a = 1$ the check computes, at four tabulated $(a_{1}, a_{2})$ and the de Sitter value, $\mathcal{R}(\gamma_{i}\wedge\gamma_{j}) = a_{1}^{2}\gamma_{i}\wedge\gamma_{j}$, $\mathcal{R}(\gamma_{0}\wedge\gamma_{i}) = a_{2}\gamma_{0}\wedge\gamma_{i}$, Weyl part zero. For a general instant this is the standard curvature of the flat Friedmann metric (cited from memory): the curvature is $(\dot a/a)^{2}$ on the space planes and $\ddot a/a$ on the time planes, linear in $(a_{1}^{2}, a_{2})$ at $a = 1$; the tabulated samples are a check of it in the store's sign, not a proof of the general formula. The Ricci map is then $\mathcal{R}(\gamma_{0}) = 3a_{2}\gamma_{0}$, $\mathcal{R}(\gamma_{i}) = (a_{2} + 2a_{1}^{2})\gamma_{i}$, and the field equations with a fluid at rest reduce to the two relations $a_{1}^{2} = (s + \Lambda)/3$ and $a_{2} = -s(1 + 3w)/6 + \Lambda/3$; the check verifies the field equations at five samples and that shifting $a_{2}$ breaks them. Substituting, $\mathcal{R} = \tfrac{\Lambda}{3}B + s\,K_{w}$ with $K_{w}$ equal to $\tfrac13$ on space planes and $-\tfrac{1+3w}{6}$ on time planes, so $\mathcal{R}$ is affine in $s$ and equals the de Sitter (or, at $\Lambda = 0$, flat) operator at $s = 0$.

Late times. The field equations' Bianchi identity gives $\dot\rho = -3\tfrac{\dot a}{a}(1+w)\rho$, so $\rho \propto a^{-3(1+w)}$. On the expanding branch $\dot a/a = ((s + \Lambda)/3)^{1/2} > 0$ never vanishes. If $\Lambda > 0$, $\dot a/a \ge (\Lambda/3)^{1/2}$ and $a \to \infty$ exponentially; if $\Lambda = 0$, $\dot a \propto a^{1 - 3(1+w)/2}$ gives $a \propto t^{2/(3(1+w))} \to \infty$. Either way $\rho$ decreases continuously to $0$ and $T \to 0$, so the local data tend to $(\tfrac{\Lambda}{3}B, 0)$, calibrated data.

Reach. Given $\beta > 0$ and $-2\beta \le \alpha < \beta$, put $w = 0$, $\Lambda = 2\alpha + \beta$ if $\alpha \ge -\beta/2$, and $w = -(2\alpha + \beta)/(3\beta)$, $\Lambda = 0$ otherwise; in both cases $w \in [0,1]$, $\Lambda \in [0, 3\beta)$, $s = 3\beta - \Lambda > 0$, and the two relations return $(a_{1}^{2}, a_{2}) = (\beta, \alpha)$ (tabulated grid in the check, with the controls $\alpha = \beta$, $\alpha < -2\beta$, $\beta \le 0$ outside the reach).

\emph{(iii).} If $f$ depends on the curvature only, each flat Friedmann field of (ii) is admissible by hypothesis and its local data tend to calibrated data, so by (i) $\eta = \tfrac14$ on it, at every instant; by the reach, $f = \tfrac14$ at every isotropic Weyl-free curvature with $\beta > 0$, $-2\beta \le \alpha < \beta$, and at $\alpha = \beta$ by the calibration itself. For the Einstein static universe the check computes the curvature from the stereographic gauge field $\hb^{-1}(e_{i}) = \tfrac{2A}{1 + |x|^{2}}e_{i}$ about $x = 0$ (a point of an open chart; the field is homogeneous, so the value holds at every point): it is Weyl-free with $\beta = 1/A^{2}$, $\alpha = 0$, equal to the flat Friedmann instant $a_{1} = 1/A$, $a_{2} = 0$, which by (ii) is the dust field with $\Lambda = 1/A^{2}$, $s = 2/A^{2}$ at its inflection instant; both satisfy the dust field equations with these values (checked at $A = 2, 3, \tfrac12$, with a mismatched radius as control). The density $\tfrac14(1 + c\,S_{ab}S^{ab})$ is a continuous function of the curvature, so it equals $\tfrac14$ there whenever the flat dust Friedmann-$\Lambda$ fields are admissible; in Theorem \claimref{synthesis.loc-independent-of-clausius-and-killing-data} its admissible class is defined by (H1) with that density, which requires $S_{ab}S^{ab}$ constant on each field, and $S_{ab}S^{ab} = \tfrac34 s^{2}$ varies along every flat dust Friedmann field, so those fields are not in it.

\emph{(iv).} Along flat dust Friedmann with $\Lambda = 0$ ($a_{1}^{2} = s/3$, $a_{2} = -s/6$) the check computes the scalar curvature and $S_{ab}S^{ab} = \mathcal{R}_{ab}\mathcal{R}^{ab} - \tfrac14\mathcal{R}^{2}$ from the gauge field at three instants: the scalar curvature takes three values and $S_{ab}S^{ab} = \tfrac34\mathcal{R}^{2}$ at each. So $\eta = \tfrac14(1 + c\,S_{ab}S^{ab}/\mathcal{R}^{2}) = \tfrac14(1 + \tfrac34 c)$ is constant on the field, and with $G_{e} = 1/(4\eta)$ the field equations hold for dust of density $\rho_{e} = \rho/G_{e}$ (the curvature fixes only $8\pi G_{e}\rho_{e}$), so the defect identity gives (H1) on every patch. At Ricci-flat data the quotient is undefined; no continuous extension takes the value $\tfrac14$ at flat data and $\tfrac14(1 + \tfrac34 c)$ along a path of data tending to flat data, by (i).

\emph{(v).} The trace-free part of $\mathcal{G}(b) + \Lambda b = -4t(b)$ is $\mathcal{R}^{\circ}(b) = -4t^{\circ}(b)$, since $\mathcal{G}$ and $\mathcal{R}$ differ by a multiple of the identity; so $X = 0$, $Q = 0$ and $\eta = \tfrac14$ on every pair satisfying the field equations at $G_{e} = 1$, and by the defect identity (H1) holds on all its patches (checked on three Friedmann samples). On the Einstein static universe of radius $2$ ($\mathcal{R}(\gamma_{0}) = 0$, $\mathcal{R}(\gamma_{i}) = \tfrac12\gamma_{i}$, scalar $\tfrac32$), with $t(b) = \sigma(b\cdot\gamma_{0})\gamma_{0}$: $\mathcal{R}^{\circ}(\gamma_{0}) = -\tfrac38\gamma_{0}$, $\mathcal{R}^{\circ}(\gamma_{i}) = \tfrac18\gamma_{i}$, $t^{\circ}(\gamma_{0}) = \tfrac34\sigma\gamma_{0}$, $t^{\circ}(\gamma_{i}) = -\tfrac14\sigma\gamma_{i}$, so $X = (\tfrac18 - \sigma)\,\mathrm{diag}(-3, 1, 1, 1)$ and $Q = 12(\tfrac18 - \sigma)^{2}$. The field equations with $G_{e} = 1/(4\eta)$, $\mathcal{G}(b) + \Lambda b = -t(b)/\eta$, give $\Lambda = \tfrac14$ and $\sigma = \eta/2$. At $\eta = \tfrac12$, $\sigma = \tfrac14$, $Q = \tfrac{3}{16}$ and $\tfrac14(1 + \tfrac{16}{3}\cdot\tfrac{3}{16}) = \tfrac12$: the density of the model takes the value $\eta$ that the field's own Clausius balance needs. The check computes $Q$, the field equations and the defect $t(k,k) + \eta\,k\cdot\mathcal{R}(k) = 0$ at five null $k$ from the gauge field, together with the second fixed point $\eta = \tfrac14$, $\sigma = \tfrac18$, $Q = 0$ (the same universe at $G_{e} = 1$), and a control (the density $\tfrac14$ with $\sigma = \tfrac14$ leaves a nonzero defect). The field is homogeneous and its symmetries fix the dust's rest frame, so $Q$ and $\eta$ are constant on it, $\eta$ is $C^{1}$, and (H1) holds on every patch on the integrated reading. $\eta$ is a polynomial in the local data, continuous everywhere, and $\tfrac14$ at calibrated data ($X = 0$ there).

\emph{(vi).} By (i) and (iii), under continuity the Clausius relation and calibration force $\eta = \tfrac14$ wherever the admissible class carries the local data to calibrated data, and by (iii) the flat Friedmann fields do so for every isotropic Weyl-free curvature in the stated range; (Loc)'s curvature clause adds content only at curvature values outside every such path. By (v), a density depending on matter is not constrained off the locus of Einstein's equations at the calibrated coupling, on any field none of whose points has data near calibrated data; (Loc)'s matter clause, read through Theorem \claimref{synthesis.einstein-equations-from-c-loc-and-clausius}, is the statement that such fields have the coupling $G_{e} = 1$ too. Neither part uses a property of the series beyond the stored theorems cited; the result is transport-absent.
